\documentclass[11pt]{article}
\usepackage[a4paper,margin=0.95in]{geometry}
\usepackage{amsmath,amssymb,amsthm,mathtools}
\usepackage{enumitem}
\usepackage{booktabs}
\usepackage{fancyhdr}

\setlist{itemsep=4pt,topsep=4pt}
\pagestyle{fancy}\fancyhf{}
\lhead{\small Parametric Inference --- Mock Semestral}
\rhead{\small Page \thepage\ of \pageref{LastPage}}
\usepackage{lastpage}
\cfoot{}

\newcommand{\E}{\mathbb{E}}
\newcommand{\Var}{\operatorname{Var}}
\newcommand{\Cov}{\operatorname{Cov}}
\newcommand{\Prob}{\mathbb{P}}
\newcommand{\R}{\mathbb{R}}
\newcommand{\Xbar}{\overline{X}}
\newcommand{\Ybar}{\overline{Y}}
\newcommand{\iid}{\overset{\text{iid}}{\sim}}
\newcommand{\ind}{\perp\!\!\!\perp}
\newcommand{\I}{\mathbb{I}}
\newcommand{\qmk}[1]{\hfill\textnormal{[\textbf{#1}]}}

\newcounter{q}
\newenvironment{question}[1]{%
  \medskip\refstepcounter{q}%
  \noindent\textbf{Q\theq.}\ \qmk{#1}\par\nopagebreak\smallskip}{\par\medskip\hrule height 0pt}

\begin{document}

\thispagestyle{empty}
\begin{center}
{\large\scshape Indian Statistical Institute}\\[10pt]
{\LARGE\bfseries Parametric Inference}\\[6pt]
{\large Mock Semestral Examination}\\[4pt]
{B.Stat.\ 3rd Year \quad$\cdot$\quad Semester 1, 2026--27}
\end{center}

\vspace{6pt}
\hrule
\vspace{6pt}

\noindent
\begin{tabular}{@{}ll@{\hspace{2.2cm}}ll@{}}
\textbf{Time:} & 3 hours & \textbf{Maximum marks:} & 100 \\
\end{tabular}

\vspace{4pt}
\hrule
\vspace{10pt}

\noindent\textbf{Instructions.}
\begin{itemize}[leftmargin=1.5em]
\item Answer \textbf{all} questions. Marks are shown in brackets against each part.
\item State clearly any theorem you invoke, and verify that its hypotheses hold. Quoting
Lehmann--Scheff\'e, Basu or Karlin--Rubin without checking completeness, ancillarity or MLR
respectively will not earn full credit.
\item Throughout, $\Phi$ and $\varphi$ denote the standard normal cdf and pdf.
\item You may use, without proof, the following values:
\[
z_{0.05}=1.645,\qquad z_{0.025}=1.960,\qquad
\Phi(0.98)=0.8365,\qquad \Phi(1.645)=0.9500 .
\]
\item Unless stated otherwise, ``risk'' means mean squared error and ``loss'' means squared
error loss.
\end{itemize}

\vspace{8pt}
\hrule
\vspace{14pt}

%%%%%%%%%%%%%%%%%%%%%%%%%%%%%%%%%%%%%% Q1
\begin{question}{12 marks}
\begin{enumerate}[label=(\alph*),leftmargin=2em]
\item State the Neyman--Fisher factorisation theorem. \qmk{3}
\item State and prove the Rao--Blackwell theorem. State precisely when the improvement is
strict. \qmk{6}
\item Give, with justification, an example of a family of distributions possessing a
sufficient statistic that is \emph{not} complete. \qmk{3}
\end{enumerate}
\end{question}

%%%%%%%%%%%%%%%%%%%%%%%%%%%%%%%%%%%%%% Q2
\begin{question}{12 marks}
Let $X\sim\mathrm{Bin}(m,\theta)$ with $m$ known and $0<\theta<1$ unknown.
\begin{enumerate}[label=(\alph*),leftmargin=2em]
\item Show that $X$ is a complete sufficient statistic for $\theta$. \qmk{4}
\item For $1\le k\le m$, find the UMVUE of $\theta^{k}$. \qmk{5}
\item Show that $\psi(\theta)=\theta^{m+1}$ admits no unbiased estimator. \qmk{3}
\end{enumerate}
\end{question}

%%%%%%%%%%%%%%%%%%%%%%%%%%%%%%%%%%%%%% Q3
\begin{question}{12 marks}
Let $X_1,\dots,X_n\iid\mathrm{Poisson}(\lambda)$, $\lambda>0$.
\begin{enumerate}[label=(\alph*),leftmargin=2em]
\item Find the UMVUE of $\Prob_\lambda\{X_1=0\}=e^{-\lambda}$. \qmk{6}
\item Compute the variance of your estimator. \qmk{3}
\item Does it attain the Cram\'er--Rao lower bound? Justify your answer, and comment on what
happens as $n\to\infty$. \qmk{3}
\end{enumerate}
\end{question}

%%%%%%%%%%%%%%%%%%%%%%%%%%%%%%%%%%%%%% Q4
\begin{question}{10 marks}
Let $X_1,\dots,X_n\iid N(\mu,\sigma^{2})$ with $n\ge3$ and both parameters unknown. Write
$S^{2}=\frac{1}{n-1}\sum_i(X_i-\Xbar)^{2}$ and $R=X_{(n)}-X_{(1)}$ for the sample range.
\begin{enumerate}[label=(\alph*),leftmargin=2em]
\item Show that $\big(\sum_iX_i,\ \sum_iX_i^{2}\big)$ is complete and sufficient. \qmk{4}
\item Show that $V=R/S$ is ancillary. \qmk{3}
\item Deduce that $\E[R]=\E[V]\,\E[S]$. \qmk{3}
\end{enumerate}
\end{question}

%%%%%%%%%%%%%%%%%%%%%%%%%%%%%%%%%%%%%% Q5
\begin{question}{12 marks}
Let $X_1,\dots,X_n\iid f(x\mid\theta)=\frac1\theta e^{-x/\theta}$, $x>0$, $\theta>0$.
Consider the inverse-gamma family of priors
\[
\pi_{\alpha,\beta}(\theta)\ \propto\ \theta^{-(\alpha+1)}e^{-\beta/\theta},
\qquad \alpha,\beta>0 .
\]
\begin{enumerate}[label=(\alph*),leftmargin=2em]
\item Show that this family is conjugate and identify the posterior. \qmk{5}
\item Find the Bayes estimate of $\theta$ under squared error loss, stating any condition
needed for it to exist. \qmk{3}
\item Is the Bayes estimate unbiased for $\theta$? \qmk{2}
\item Find the generalized Bayes estimate of $\theta$ under the improper prior
$\pi(\theta)=1/\theta$, and comment on how it compares with $\Xbar$. \qmk{2}
\end{enumerate}
\end{question}

%%%%%%%%%%%%%%%%%%%%%%%%%%%%%%%%%%%%%% Q6
\begin{question}{12 marks}
Let $X\sim\mathrm{Bin}(n,\theta)$, $0<\theta<1$, and let the prior be
$\mathrm{Beta}(p,q)$.
\begin{enumerate}[label=(\alph*),leftmargin=2em]
\item Show that the Bayes estimate of $\theta$ under squared error loss is
$T_{p,q}=\dfrac{p+X}{p+q+n}$. \qmk{3}
\item Compute $R(\theta,T_{p,q})$ and determine all $(p,q)$ for which the risk is constant in
$\theta$. \qmk{6}
\item Hence identify a minimax estimator of $\theta$, and compare its maximum risk with that
of the UMVUE $X/n$. \qmk{3}
\end{enumerate}
\end{question}

%%%%%%%%%%%%%%%%%%%%%%%%%%%%%%%%%%%%%% Q7
\begin{question}{12 marks}
A single observation $X$ is available. Consider
\[
H_0:\ X\sim N(0,1)\qquad\text{against}\qquad H_A:\ X\sim N(0,4).
\]
\begin{enumerate}[label=(\alph*),leftmargin=2em]
\item Derive the form of the most powerful test of level $\alpha$. \qmk{6}
\item For $\alpha=0.05$, state the rejection region explicitly and compute the power.
\qmk{3}
\item Now let $X\sim N(0,\sigma^{2})$ with $\sigma^{2}>0$ unknown. Does a UMP test of size
$\alpha$ exist for $H_0:\sigma^{2}=1$ against $H_A:\sigma^{2}\ne1$? Justify carefully.
\qmk{3}
\end{enumerate}
\end{question}

%%%%%%%%%%%%%%%%%%%%%%%%%%%%%%%%%%%%%% Q8
\begin{question}{10 marks}
Let $X_1,\dots,X_n\iid \mathrm{Uniform}[0,\theta]$, $\theta>0$.
\begin{enumerate}[label=(\alph*),leftmargin=2em]
\item Show that the family has monotone likelihood ratio in $X_{(n)}$. \qmk{4}
\item Derive the UMP test of size $\alpha$ for $H_0:\theta\le\theta_0$ against
$H_A:\theta>\theta_0$. \qmk{4}
\item Write down its power function and verify that the size is attained at
$\theta=\theta_0$. \qmk{2}
\end{enumerate}
\end{question}

%%%%%%%%%%%%%%%%%%%%%%%%%%%%%%%%%%%%%% Q9
\begin{question}{10 marks}
\begin{enumerate}[label=(\alph*),leftmargin=2em]
\item Let $X_1,\dots,X_n\iid\mathrm{Cauchy}(\theta)$, i.e.\ with density
$\frac{1}{\pi\{1+(x-\theta)^{2}\}}$. Show that $\Xbar$ is \emph{not} a consistent estimator
of $\theta$, whereas the sample median is. \qmk{6}
\item In the model $Y_i=\alpha+\beta x_i+\varepsilon_i$, $i=1,\dots,n$, with $x_i$ known
constants and $\varepsilon_i$ iid with mean $0$ and variance $\sigma^{2}$, state and justify
the condition on $\{x_i\}$ under which the least squares estimate $\hat\beta_{LS}$ is
consistent for $\beta$. \qmk{4}
\end{enumerate}
\end{question}

\vspace{20pt}
\begin{center}
\rule{0.35\textwidth}{0.4pt}\\[4pt]
\textit{End of paper}
\end{center}

\end{document}
